\section{The Propers: Themes and Movement}

Christ is King of all by nature, as the eternal Son, and as man he receives
the same dominion over the nations, asked for and won by his cross. His
kingdom takes neither its origin nor its means from this world, and so it
reaches every nation and every ruler; when he hands it over to the Father it
does not end, because what he hands over is the people he has redeemed,
brought to the vision of God.

\textbf{The slain Lamb receives what the Son always had.} The Introit takes up
the song of the angels round the throne when the Lamb has taken the sealed
book that no one else was worthy to open: \latin{Dignus est Agnus, qui occisus
est}, worthy \latin{accipere \dots\ divinitatem}, to receive divinity (Apoc
5:12). The Offertory gives the Father's word in the second psalm, ``Ask of me,
and I will give thee the Gentiles for thy inheritance,'' spoken to the king
who has just recited the decree, ``Thou art my son, this day have I begotten
thee'' (Ps 2:7--8). Of the eternal Son, both seem to say that he receives what
he could never lack. The Alleluia sings only the second half of Dan 7:14, whose
first half begins ``And he gave him power.''

The Fathers answer from the psalm, and they divide its two verses
differently. St Hilary of Poitiers, on St Paul's word at Antioch (Acts 13:33),
takes the ``today'' of the decree as the begetting of the firstborn from the
dead, and what the risen Son asks and is given is \latin{ut esset, quod fuerat
ante}, that he might be what he had been before (\work{Tractatus super
Psalmos} 2.30--33). St Augustine reads verse~7 of the eternal generation and
only verse~8 in time, of ``the Manhood which He took on Himself, who offered up
Himself as a Sacrifice in the stead of all sacrifices'' (\work{Enarrationes in
Psalmos} 2.6--7). St Thomas Aquinas states the division as a principle:
Christ holds the nations \latin{duplici jure: scilicet hereditario, et
meritorio}, by inheritance, which is the eternal generation, and by merit,
which is his asking in prayer and in his passion; God willed that he should
ask \latin{quod sibi jure hereditario competebat}, what was already his by
inheritance (\work{Postilla super Psalmos} 2, nn.~5--6). St Bede has the Son of man receive the dispensation of the Incarnation
that he himself, as God, had disposed with the Father: the Lion and the Lamb
are one Lord (\work{Expl.\ Apoc.}\ I, on 5:5--7).

Taken with Augustine's words on verse~8, the Secret offers the asking of the
psalm: the \latin{hostiam \dots\ humanae reconciliationis} is the Son himself,
\latin{quem sacrificiis praesentibus immolamus}, and the \latin{unitatis et
pacis dona} asked of him for all nations are the inheritance Hilary describes,
the nations carried over out of the unjust right of those who lorded it over
them into an eternal and divine kingdom (2.31). His words for that
inheritance are the Epistle's: those translated
``into the kingdom of the Son of his love'' by ``the firstborn from the
dead,'' who makes ``peace through the blood of his cross'' (Col 1:13, 18, 20).

\textbf{The kingdom is not from hence, and is therefore everywhere.} In the
governor's hall, bound and delivered up by his own nation, Jesus says that his
kingdom is not \latin{de hoc mundo}; that were it \latin{ex hoc mundo}, his
servants would strive for him; and that now it is not \latin{hinc}, ``from
hence'' (Jn 18:36). The Gradual before the Gospel gives him dominion ``from sea
to sea'' and sings that ``all kings of the earth shall adore him: all nations
shall serve him'' (Ps 71:8, 11). The Collect asks that the families of the
nations, \latin{peccati vulnere disgregatae}, scattered by the wound of sin, be
subjected to his \latin{suavissimo imperio}, his most gentle rule. The
Postcommunion has the faithful \latin{sub Christi Regis vexillis militare},
fighting under the standards of Christ the King, whose servants, the Gospel
says, do not fight.

Augustine draws the answer from the Lord's exact words. He did not say that
his kingdom is not \emph{in} this world, but not \emph{of} it; not that it is
not \emph{here}, but not \emph{from hence}. ``For His kingdom is here until the
end of the world, having tares intermingled therewith until the harvest,'' and
yet ``it only sojourns as a stranger in the world'' (\work{Tractates on the
Gospel of John} 115.2). The kingdom is ``Those who believe in Him,'' carried
over, in the Epistle's verse that Augustine sets beside the Gospel, from the
power of darkness into the kingdom of the Son; its subjects are those who are
of the truth and hear his voice, where hearing is obeying and believing
(115.4). St John Chrysostom takes the kingdom to be the royal power rather
than the people: Christ rules here, but ``He has his empire from above,'' and
the servants who do not fight are the proof, for earthly kingship's ``strength
lies in servants; but that which is above is sufficient for itself, needing
nothing'' (\work{Homilies on John} 83.4). Aquinas keeps both and concludes
\latin{non est hinc \dots\ est tamen hic, quia ubique est}: not from here,
and yet here, because everywhere, with the Offertory's verse in support
(\work{Super Io.}\ 18 lect.~6).

The kings of the chants have a present office as well as a final one.
Augustine asks why the kings of Tharsis and Saba ``shall lead'' their gifts, and
answers that the gifts ``seem to me to signify
men, whom into the fellowship of the Church of Christ the authority of kings
doth lead''; the kings who adore belong to the poor people who confess their
sins (\work{Enarrationes in Psalmos} 71.13--14). In \work{Letter} 185, written
to justify the imperial laws against the Donatists, he sets the second
psalm's command to the kings beside the Gradual's verse: ``a man serves God in
one way in that he is man, in another way in that he is also king''
(185.19--20). St Robert Bellarmine finds the adoration of all kings fulfilled
``to a certain extent'' as princes and nations are converted, ``to the letter''
at the day of judgment, and ``always'' in ``the actual power that Christ has
over all princes and all nations'' (\work{Explanatio}, Ps~71). In
his note on the Postcommunion, Blessed Ildefonso Schuster names the standard:
``The Royal Standard of Jesus Christ is the Cross'' (p.~204). By Chrysostom's
argument the faithful's warfare cannot be the King's defence; by Hilary's,
which places the kingdom of God where sin is overcome and death destroyed
(\work{Tractatus super Psalmos} 2.42), it is a share in his victory over
sin.

\textbf{The kingdom handed over to the Father has no end.} The Preface gives
thanks that the Father anointed his Son \latin{Sacerdotem aeternum et
universorum Regem}, so that, offering himself on the altar of the cross, he
might accomplish man's redemption and then, \latin{suo subiectis imperio
omnibus creaturis}, hand over to the Father ``an eternal and universal
kingdom, a kingdom of truth and life; a kingdom of holiness and grace, a
kingdom of justice, love and peace.'' Its \latin{traderet} is the ``delivered
up'' of 1 Cor 15:24. Against it stand the Alleluia's \latin{potestas aeterna,
quae non auferetur} and the Communion's \latin{Sedebit Dominus Rex in
aeternum}. St Cyril of Jerusalem calls the teaching that Christ's kingdom shall
end ``another head of the dragon,'' and reads the Apostle's ``until'' (1 Cor
15:25) as setting no term, as death reigned from Adam ``until'' Moses (Rom 5:14)
(\work{Catechetical Lectures} 15.27--32).

Augustine answers by asking what the kingdom is. To deliver it up is ``as much
as to say, When He shall have brought believers to the contemplation of God,
even the Father'' (\work{De Trinitate} I.8.16), for ``His faithful ones, whom He
has redeemed with His blood, are called His kingdom'' (I.10.21). Hilary
reaches the same answer from the wording of the verse: ``It is not said, He
shall deliver up His Kingdom, but, He shall deliver up the Kingdom, that is,
deliver up to God us who have been made the Kingdom by the glorifying of His
body'' (\work{De Trinitate} XI.39). Both cite ``Blessed are the pure in heart,
for they shall see God'' for the purity that vision requires. They part on the
Son's subjection: Augustine keeps the form of the servant, the Priest, as what
is subject, while Hilary says the Son ``must now become God only; not,
however, by casting off the body, but by translating it through subjection''
(XI.40). Set beside Hilary's tractate on the second psalm, where the nations
given to the Son are the inheritance of the firstborn from the dead, his
\work{De Trinitate} makes what is given to the Son what the Son gives to the
Father: the Offertory's inheritance is the kingdom the Preface has him hand
over.

Psalm~28 ends, ``The Lord maketh the flood to dwell: and the Lord shall sit
king for ever''; the Communion keeps the session and leaves out the flood, and
its expositors place the session by what they take the flood to be. Augustine: ``The
Lord therefore first inhabiteth the deluge of this world in His Saints, kept
safely in the Church, as in the ark,'' and ``afterward He will sit reigning in
them for ever'' (\work{Enarrationes in Psalmos} 28.10). St Basil calls the
grace of baptism the deluge, and God, settled in the soul that shines from it,
makes it, as it were, a throne for himself (\work{Homilia in Psalmum} 28.8).
Bellarmine, on his second interpretation, places the session after the
judgment, ``all power in general being swept away'' (\work{Explanatio},
Ps~28). Augustine's session runs on from one to the other;
Bellarmine's begins when the judgment is done.

The first two answers meet at Jn 18:37. ``For this was I born''
Augustine reads of the birth from the Virgin (115.4), Chrysostom of the
eternal generation: the Son, ``born a king,'' ``has nothing which He received
in addition,'' and every text of giving refers to his generation. Augustine's
division of the psalm grants that premise and refers the receiving to the
humanity; in his two expositions the eternal birth grounds the decree of Ps
2:7, and the temporal birth the testimony before Pilate.
